\section{Diseño e implementación en SP1}
\label{sec:sp1_impl}

La implementación en SP1 expresa los dos enunciados presentados en la
Sección~\ref{sec:puente_arquitectura} mediante dos programas invitados escritos en Rust, cuya
ejecución se demuestra en la máquina virtual descrita en la Sección~\ref{sec:sp1}, y dos programas
anfitriones que preparan sus entradas y solicitan las pruebas. La arquitectura, los contratos y el
flujo de una transferencia son los descritos en la Sección~\ref{sec:lock_and_mint}, por lo que esta sección
sigue la estructura de la Sección~\ref{sec:noir_impl} y se ocupa de lo que demuestra cada programa,
de la codificación y la verificación de sus valores públicos y de la generación de las pruebas. La
descripción corresponde al código del repositorio del puente y de los repositorios de sus dos
programas~\cite{zkguaranibridgesp1, sp1synccommitteerepo, txinclusionsp1repo}. En ninguno de los
dos programas las dependencias criptográficas, que son las bibliotecas \texttt{sha2},
\texttt{sha3} y \texttt{bls12\_381\_plus}, se sustituyen por las versiones modificadas que dirigen
sus operaciones hacia los precompilados de la máquina virtual
(Sección~\ref{sec:sp1_flujo}), de modo que el cálculo de SHA-256, de Keccak-256 y de las
operaciones sobre BLS12-381 se ejecuta como instrucciones RISC-V ordinarias.

\subsection{Programa de consenso}
\label{sec:sp1_impl_consenso}

El programa invitado de consenso corresponde al paquete \texttt{sync-committee-program}, y su
anfitrión es el ejecutable \texttt{sync-committee}, que lee un archivo con los datos de la cadena
de baliza y escribe en la entrada del invitado doce valores. Esos valores son los cinco campos del
encabezado atestiguado, el dominio de firma, la firma agregada, las claves públicas comprimidas
de los miembros del comité, el vector de participación, el hash del bloque de ejecución que se
pretende anclar, la rama SSZ que lo vincula con \texttt{body\_root} y su índice
generalizado~\cite{sp1synccommitteerepo}. Cuando la cadena de origen dispone de un nodo de baliza,
el puente obtiene esos datos tomando como encabezado atestiguado el del bloque de baliza cuyo
\textit{execution payload} contiene el hash del bloque del bloqueo, y tomando la firma y el vector
de participación del bloque siguiente, ya que el agregado incluido en un \textit{slot} firma la
raíz del bloque anterior~\cite{zkguaranibridgesp1, ethconsensusaltair}. El anfitrión resuelve el
dominio, a partir de un valor explícito o de la versión de la bifurcación y de la raíz de los
validadores del bloque génesis, y verifica de forma nativa la rama y la firma antes de solicitar
la ejecución, con el fin de rechazar entradas inconsistentes sin esperar a la máquina virtual.
Estas comprobaciones del anfitrión quedan fuera de la ejecución demostrada
(Sección~\ref{sec:sp1_programas}), por lo que el invitado repite todas las que sustentan el
enunciado, con la única excepción de la procedencia del dominio, que se contrasta en el contrato.

El invitado calcula la raíz del encabezado mediante \texttt{hash\_tree\_root} y obtiene la raíz de
firma como el hash SHA-256 de la concatenación de esa raíz con el dominio~\cite{ethconsensusphase0}.
Suma después las claves de los miembros cuyo bit de participación vale uno, descomprimiendo
únicamente las de esos miembros, y verifica la Ecuación~\eqref{eq:bls_agregada} sobre la raíz de
firma, con la función $H_{\mathbb{G}_2}$ implementada como hash a la curva con la etiqueta de
separación de dominio que emplea Ethereum~\cite{rfc9380}. A diferencia de un circuito, la
ejecución solo recorre las claves de los firmantes, de manera que el trabajo de esta etapa depende
de la participación efectiva. Verifica también, mediante la función \texttt{verify\_merkle\_branch},
que la rama conduce desde el hash del bloque de ejecución hasta \texttt{body\_root} y que el índice
generalizado corresponde exactamente a la profundidad de la rama, y comprueba que el número de
firmantes multiplicado por tres supere el doble del tamaño del comité. Si las claves o el vector de
participación están mal formados, el invitado no interrumpe su ejecución, sino que registra cero
firmantes, de modo que la prueba se genera igualmente con un resultado negativo. El invitado
compromete diez valores públicos, entre ellos el indicador \texttt{valid}, que solo es verdadero si
la firma, la rama y el umbral se verifican, y \texttt{committeeCommitment}, el hash SHA-256 de las
claves comprimidas del comité concatenadas en orden~\cite{sp1synccommitteerepo}.

\texttt{BeaconLightClient} verifica la prueba con la clave del programa de consenso, exige que
\texttt{valid} sea verdadero, compara el dominio, el índice generalizado y el compromiso del comité
con los valores fijados por su propietario, rechaza un hash de bloque o un \textit{slot} nulos y, si
todas las comprobaciones se cumplen, registra el \textit{slot} asociado al hash del bloque de
ejecución~\cite{zkguaranibridgesp1}. El encabezado que firma el comité es el atestiguado. Ni el
invitado ni el contrato verifican la rama que vincularía un encabezado finalizado con el
\texttt{state\_root} del atestiguado (Sección~\ref{sec:eth_cadena_verificacion}), de modo que el
anclaje acredita la firma del comité sobre el bloque, pero no su finalidad.

\subsection{Programa de inclusión y ejecución exitosa}
\label{sec:sp1_impl_inclusion}

\begin{sloppypar}
El programa invitado de inclusión corresponde al paquete \texttt{tx-inclusion-verifier-program}, cuya
función principal lee una estructura de entrada, delega la verificación en la función
\texttt{process\_input} de la biblioteca \texttt{tx-inclusion-verifier-lib} y compromete el
resultado~\cite{txinclusionsp1repo}. Su anfitrión, el ejecutable \texttt{tx-inclusion-verifier},
recibe esa entrada en un archivo que prepara el \textit{relayer}. Este reconstruye el encabezado
del bloque y comprueba que su hash Keccak-256 coincida con el que informa el nodo, reserializa cada
transacción y comprueba su hash, obtiene los recibos del bloque, comprueba que cada uno corresponda
a la transacción de su posición y los serializa en la codificación que almacena el árbol de
recibos, con el byte de tipo antepuesto en los recibos tipados~\cite{zkguaranibridgesp1, eip2718}.
\end{sloppypar}

El invitado calcula el hash Keccak-256 del encabezado, que identifica el bloque, y extrae de ese
mismo encabezado \texttt{transactionsRoot} y \texttt{receiptsRoot} (Tabla~\ref{tab:eth_header}). Con la lista completa de transacciones reconstruye el
árbol de transacciones, cuyas claves son las codificaciones RLP de las posiciones, comprueba que su
raíz coincida con la del encabezado y que la transacción declarada ocupe la posición indicada, y
repite el procedimiento con los recibos y su raíz. El recibo que examina es el de esa misma
posición, por lo que presentar el recibo exitoso de otra transacción alteraría la raíz reconstruida
o correspondería a otra posición. La biblioteca admite también la verificación de una rama del
árbol, pero el puente utiliza la reconstrucción completa. El invitado decodifica después el sobre
de la transacción, del que obtiene el destinatario, el valor y los datos de la llamada, cuyo hash y
longitud compromete, y lee el código de estado del recibo, que desde Byzantium indica si la
ejecución terminó sin revertir~\cite{eip658}. Un recibo anterior a esa actualización, que contiene
una raíz de estado en lugar del código, se considera no válido~\cite{txinclusionsp1repo}.

El estado exitoso no se exige mediante una aserción del invitado. La función
\texttt{process\_input} no interrumpe la ejecución ante un error, sino que lo refleja en los
indicadores \texttt{receiptVerified} y \texttt{txSucceeded}, y este último solo es verdadero si el
recibo se verificó y su código indica éxito. Una prueba de un bloqueo revertido es, por lo tanto,
una prueba válida, y es \texttt{ZKReceiver} quien la rechaza con el error
\texttt{TransactionFailed}~\cite{zkguaranibridgesp1}. El \textit{relayer}, por su parte, descarta un
bloqueo revertido antes de generar las pruebas y vuelve a comprobar ambos indicadores antes de
enviarlas, pero esa conducta solo evita trabajo innecesario. La condición que impide acuñar es la
del receptor, que se aplica a cualquier parte que presente la prueba. Los datos de la llamada se
vinculan con el destinatario y el monto del mismo modo descrito en la
Sección~\ref{sec:puente_aceptacion}, ya que el receptor recibe esos datos, comprueba su longitud y
su hash frente a los comprometidos y decodifica de ellos la llamada a \texttt{lock}. El programa
no examina los eventos del recibo, a diferencia del circuito de inclusión de Noir
(Sección~\ref{sec:noir_impl_inclusion}). El evento \texttt{Locked} sirve al \textit{relayer} para
detectar la operación, mientras que el receptor se apoya en el recibo probado y en los datos de la
llamada.

\subsection{Valores públicos y verificación en Solidity}
\label{sec:sp1_impl_publicas}

Los valores públicos de cada programa son una secuencia de bytes con la codificación del ABI de
Ethereum. El invitado de consenso construye una estructura declarada con la biblioteca
\texttt{alloy-sol-types}, cuyo cuarto campo es \texttt{executionBlockHash}, la codifica y la
compromete mediante \lstinline|commit_slice| (Sección~\ref{sec:sp1_programas}). Como todos los
campos son de tamaño fijo, la codificación ocupa exactamente diez palabras de 32 bytes y coincide
byte a byte con la que decodifica \lstinline|abi.decode| en Solidity, de modo que el hash del bloque
ocupa la cuarta palabra. El anfitrión escribe esos bytes junto con la prueba en un archivo, y el
\textit{relayer} los entrega como argumento de \lstinline|mintRemoteWithHeader|.
\texttt{ZKReceiver} los reenvía a \lstinline|updateHeader|, donde \texttt{BeaconLightClient}
verifica la prueba, decodifica la estructura y registra el anclaje con ese hash como clave. Los
valores del programa de inclusión ocupan quince palabras y llevan el hash del bloque en la quinta.
\texttt{ZKReceiver} los decodifica y consulta \lstinline|isAnchored| con ese valor, de modo que el
vínculo entre ambas pruebas es la igualdad de esos 32 bytes, comprobada a través del registro de
anclajes~\cite{zkguaranibridgesp1, sp1synccommitteerepo, txinclusionsp1repo}.

Ambos contratos delegan la verificación en la función \lstinline|verifyProof| del contrato
verificador, que recibe la clave del programa, los valores públicos y la prueba, y revierte si la
prueba no es válida para esa combinación~\cite{sp1contracts}. Las claves de los dos programas se
fijan como valores inmutables al desplegar los contratos, por lo que una prueba válida de otro
programa con valores de la misma forma se rechaza. El verificador incorpora los valores públicos a
través de su hash SHA-256 (Sección~\ref{sec:sp1_flujo}), de manera que alterar un solo byte de la
secuencia invalida la prueba. La Tabla~\ref{tab:valores_publicos} enumera los campos de ambos
programas y el uso que hacen de ellos los contratos. Algunos no intervienen en ninguna
comprobación y se exponen con fines de trazabilidad, como las raíces que
\texttt{BeaconLightClient} registra en el evento de anclaje o el hash de la transacción que
\texttt{ZKReceiver} registra en el evento \texttt{Minted}.

\begin{table}[h!]
	\centering
	\caption{Valores públicos de los programas de consenso y de inclusión de SP1 y uso que hacen de
	ellos los contratos de destino.}
	\label{tab:valores_publicos}
	\footnotesize
	\begin{tabular}{>{\raggedright\arraybackslash}p{3.6cm} >{\raggedright\arraybackslash}p{5.9cm}
		>{\raggedright\arraybackslash}p{6.0cm}}
		\toprule
		\textbf{Valor} & \textbf{Contenido} & \textbf{Uso en el contrato} \\
		\midrule
		\multicolumn{3}{l}{\textit{Programa de consenso, consumido por} \texttt{BeaconLightClient}} \\
		\addlinespace
		\texttt{valid} & Conjunción de la verificación de la firma, de la rama SSZ y del umbral de
		participación. & Debe ser verdadero. \\
		\texttt{signingRoot} & Mensaje firmado por el comité. & Se registra en el evento de anclaje. \\
		\texttt{headerRoot} & Raíz del encabezado de la cadena de baliza. & Se registra en el evento
		de anclaje. \\
		\texttt{executionBlockHash} & Hash del bloque de ejecución. & Se ancla y no puede ser cero. \\
		\texttt{domain} & Dominio de firma. & Debe coincidir con el fijado por el propietario. \\
		\texttt{slot} & \textit{Slot} del encabezado. & Se registra con el anclaje y no puede ser
		cero. \\
		\texttt{blockHashGindex} & Índice generalizado de la rama SSZ. & Debe coincidir con el fijado
		por el propietario. \\
		\texttt{committeeCommitment} & Hash SHA-256 de las claves del comité. & Debe coincidir con el
		fijado por el propietario. \\
		\texttt{committeeSize}, \texttt{participants} & Tamaño del comité y número de firmantes. & No
		se comprueban, porque el umbral forma parte de \texttt{valid}. \\
		\midrule
		\multicolumn{3}{l}{\textit{Programa de inclusión, consumido por} \texttt{ZKReceiver}} \\
		\addlinespace
		\texttt{verified} & Resultado de la comprobación de inclusión. & Debe ser verdadero. \\
		\texttt{headerBound} & La raíz de transacciones proviene del encabezado. & Debe ser
		verdadero. \\
		\texttt{txParsed} & El sobre de la transacción pudo decodificarse. & Debe ser verdadero. \\
		\texttt{txType} & Tipo de la transacción. & No se comprueba. \\
		\texttt{blockHash} & Hash Keccak-256 del encabezado de ejecución. & Debe estar anclado y forma
		parte de la clave contra la reutilización. \\
		\texttt{transactionsRoot} & Raíz del árbol de transacciones. & El programa la contrasta con
		la raíz reconstruida, y el contrato no la compara por separado. \\
		\texttt{txIndex} & Posición de la transacción en el bloque. & Forma parte de la clave contra la
		reutilización. \\
		\texttt{txHash} & Hash Keccak-256 de la transacción. & Se registra en el evento
		\texttt{Minted}. \\
		\texttt{txTo} & Destinatario de la transacción. & Debe ser el contrato \texttt{Sender}. \\
		\texttt{txValue} & Valor transferido en la moneda nativa. & No se comprueba. \\
		\texttt{txDataHash} & Hash Keccak-256 de los datos de la llamada. & Debe coincidir con el de los
		datos recibidos. \\
		\texttt{txDataLen} & Longitud de los datos de la llamada. & Debe coincidir con la de los datos
		recibidos. \\
		\texttt{receiptVerified} & El recibo de la misma posición pertenece al árbol de recibos del
		encabezado. & Debe ser verdadero. \\
		\texttt{txSucceeded} & Ese recibo registra una ejecución exitosa. & Debe ser verdadero. \\
		\texttt{receiptsRoot} & Raíz del árbol de recibos. & El programa la contrasta con la raíz
		reconstruida, y el contrato no la compara por separado. \\
		\bottomrule
	\end{tabular}
\end{table}

El indicador \texttt{txSucceeded} implica \texttt{receiptVerified},
porque el programa solo lo activa sobre un recibo verificado, y \texttt{ZKReceiver} comprueba
ambos, de modo que un recibo no verificado se rechaza con \texttt{ReceiptNotVerified} antes de
examinar el éxito. Las raíces \texttt{transactionsRoot} y \texttt{receiptsRoot} se contrastan dentro
del programa con las reconstruidas, y ningún contrato las compara por separado.

\subsection{Generación de pruebas e integración}
\label{sec:sp1_impl_integracion}

La preparación se realiza una vez por programa. La compilación de cada anfitrión compila también el
invitado para la arquitectura RISC-V de la máquina virtual, y la clave de verificación se deriva de
forma determinista del ejecutable resultante (Sección~\ref{sec:sp1_programas}). El script de
despliegue de la cadena de destino lee esas claves de los ejecutables y las fija en
\texttt{BeaconLightClient} y \texttt{ZKReceiver}, junto con la dirección del contrato verificador
de SP1 de la red y los parámetros de política. Los scripts del puente recompilan los programas
cuando sus fuentes son más recientes que los ejecutables, y cualquier modificación de un invitado,
como la incorporación de la verificación del recibo, produce una clave distinta y exige desplegar
de nuevo los contratos que la fijan~\cite{zkguaranibridgesp1}.

Durante cada transferencia, el \textit{relayer} prepara las dos entradas y ejecuta ambos
anfitriones en el modo destinado a Ethereum, que genera la prueba de la ejecución y la envuelve en
una prueba Groth16 sobre la curva BN254 (Sección~\ref{sec:sp1_flujo}). Los anfitriones habilitan
la característica \texttt{native-gnark} del kit, que permite realizar la envoltura sin un
contenedor, y verifican localmente cada prueba antes de escribirla. El \textit{relayer} genera
siempre ambas pruebas y consulta el anclaje del bloque solo después de obtenerlas. Si el bloque no
está anclado, envía las dos a \lstinline|mintRemoteWithHeader|, que ancla el bloque y acuña en una
sola transacción, mientras que, si ya lo está, envía solo la prueba de inclusión a
\lstinline|mintRemote|, y la de consenso, aunque se generó, queda sin utilizar. En ambos casos acompaña los datos de la llamada a \texttt{lock}. Los anfitriones
ofrecen, además, un modo que ejecuta el invitado sin generar ninguna prueba y otro que genera la
prueba comprimida, que no está destinada a la verificación en Ethereum.

Estos modos, junto con las transacciones de la cadena de destino, permiten distinguir la ejecución
sin prueba, que informa el número de ciclos, la generación de la prueba comprimida, el tiempo
adicional de la envoltura Groth16, que se estima por diferencia entre los tiempos de los modos
destinado a Ethereum y comprimido, $T_{\mathrm{Ethereum}} - T_{\mathrm{comprimido}}$, la
verificación local y la verificación en la cadena de destino. Los anfitriones no cronometran
internamente sus etapas, y en los dos modos que generan pruebas ejecutan en secuencia la
preparación de las claves, la generación de la prueba, su verificación local y la escritura de los
archivos de salida~\cite{sp1synccommitteerepo, txinclusionsp1repo}. Si se cronometra cada modo
completo, la diferencia es, por lo tanto, una estimación, que incluye además la diferencia entre el
costo de verificar localmente una prueba Groth16 y el de verificar una prueba comprimida. Del mismo
modo, el tiempo total del \textit{relayer} incluye todo el trabajo que efectivamente realiza, aun
cuando una de las pruebas no se envíe, mientras que los tiempos de generación de las pruebas de
consenso y de inclusión se comparan por separado. La sección siguiente define cómo se miden estas
etapas y su correspondencia con las de la implementación en Noir.
